\documentclass[11pt]{article}

    \usepackage[breakable]{tcolorbox}
    \usepackage{parskip} % Stop auto-indenting (to mimic markdown behaviour)
    

    % Basic figure setup, for now with no caption control since it's done
    % automatically by Pandoc (which extracts ![](path) syntax from Markdown).
    \usepackage{graphicx}
    % Keep aspect ratio if custom image width or height is specified
    \setkeys{Gin}{keepaspectratio}
    % Maintain compatibility with old templates. Remove in nbconvert 6.0
    \let\Oldincludegraphics\includegraphics
    % Ensure that by default, figures have no caption (until we provide a
    % proper Figure object with a Caption API and a way to capture that
    % in the conversion process - todo).
    \usepackage{caption}
    \DeclareCaptionFormat{nocaption}{}
    \captionsetup{format=nocaption,aboveskip=0pt,belowskip=0pt}

    \usepackage{float}
    \floatplacement{figure}{H} % forces figures to be placed at the correct location
    \usepackage{xcolor} % Allow colors to be defined
    \usepackage{enumerate} % Needed for markdown enumerations to work
    \usepackage{geometry} % Used to adjust the document margins
    \usepackage{amsmath} % Equations
    \usepackage{amssymb} % Equations
    \usepackage{textcomp} % defines textquotesingle
    % Hack from http://tex.stackexchange.com/a/47451/13684:
    \AtBeginDocument{%
        \def\PYZsq{\textquotesingle}% Upright quotes in Pygmentized code
    }
    \usepackage{upquote} % Upright quotes for verbatim code
    \usepackage{eurosym} % defines \euro

    \usepackage{iftex}
    \ifPDFTeX
        \usepackage[T1]{fontenc}
        \IfFileExists{alphabeta.sty}{
              \usepackage{alphabeta}
          }{
              \usepackage[mathletters]{ucs}
              \usepackage[utf8x]{inputenc}
          }
    \else
        \usepackage{fontspec}
        \usepackage{unicode-math}
    \fi

    \usepackage{fancyvrb} % verbatim replacement that allows latex
    \usepackage{grffile} % extends the file name processing of package graphics
                         % to support a larger range
    \makeatletter % fix for old versions of grffile with XeLaTeX
    \@ifpackagelater{grffile}{2019/11/01}
    {
      % Do nothing on new versions
    }
    {
      \def\Gread@@xetex#1{%
        \IfFileExists{"\Gin@base".bb}%
        {\Gread@eps{\Gin@base.bb}}%
        {\Gread@@xetex@aux#1}%
      }
    }
    \makeatother
    \usepackage[Export]{adjustbox} % Used to constrain images to a maximum size
    \adjustboxset{max size={0.9\linewidth}{0.9\paperheight}}

    % The hyperref package gives us a pdf with properly built
    % internal navigation ('pdf bookmarks' for the table of contents,
    % internal cross-reference links, web links for URLs, etc.)
    \usepackage{hyperref}
    % The default LaTeX title has an obnoxious amount of whitespace. By default,
    % titling removes some of it. It also provides customization options.
    \usepackage{titling}
    \usepackage{longtable} % longtable support required by pandoc >1.10
    \usepackage{booktabs}  % table support for pandoc > 1.12.2
    \usepackage{array}     % table support for pandoc >= 2.11.3
    \usepackage{calc}      % table minipage width calculation for pandoc >= 2.11.1
    \usepackage[inline]{enumitem} % IRkernel/repr support (it uses the enumerate* environment)
    \usepackage[normalem]{ulem} % ulem is needed to support strikethroughs (\sout)
                                % normalem makes italics be italics, not underlines
    \usepackage{soul}      % strikethrough (\st) support for pandoc >= 3.0.0
    \usepackage{mathrsfs}
    

    
    % Colors for the hyperref package
    \definecolor{urlcolor}{rgb}{0,.145,.698}
    \definecolor{linkcolor}{rgb}{.71,0.21,0.01}
    \definecolor{citecolor}{rgb}{.12,.54,.11}

    % ANSI colors
    \definecolor{ansi-black}{HTML}{3E424D}
    \definecolor{ansi-black-intense}{HTML}{282C36}
    \definecolor{ansi-red}{HTML}{E75C58}
    \definecolor{ansi-red-intense}{HTML}{B22B31}
    \definecolor{ansi-green}{HTML}{00A250}
    \definecolor{ansi-green-intense}{HTML}{007427}
    \definecolor{ansi-yellow}{HTML}{DDB62B}
    \definecolor{ansi-yellow-intense}{HTML}{B27D12}
    \definecolor{ansi-blue}{HTML}{208FFB}
    \definecolor{ansi-blue-intense}{HTML}{0065CA}
    \definecolor{ansi-magenta}{HTML}{D160C4}
    \definecolor{ansi-magenta-intense}{HTML}{A03196}
    \definecolor{ansi-cyan}{HTML}{60C6C8}
    \definecolor{ansi-cyan-intense}{HTML}{258F8F}
    \definecolor{ansi-white}{HTML}{C5C1B4}
    \definecolor{ansi-white-intense}{HTML}{A1A6B2}
    \definecolor{ansi-default-inverse-fg}{HTML}{FFFFFF}
    \definecolor{ansi-default-inverse-bg}{HTML}{000000}

    % common color for the border for error outputs.
    \definecolor{outerrorbackground}{HTML}{FFDFDF}

    % commands and environments needed by pandoc snippets
    % extracted from the output of `pandoc -s`
    \providecommand{\tightlist}{%
      \setlength{\itemsep}{0pt}\setlength{\parskip}{0pt}}
    \DefineVerbatimEnvironment{Highlighting}{Verbatim}{commandchars=\\\{\}}
    % Add ',fontsize=\small' for more characters per line
    \newenvironment{Shaded}{}{}
    \newcommand{\KeywordTok}[1]{\textcolor[rgb]{0.00,0.44,0.13}{\textbf{{#1}}}}
    \newcommand{\DataTypeTok}[1]{\textcolor[rgb]{0.56,0.13,0.00}{{#1}}}
    \newcommand{\DecValTok}[1]{\textcolor[rgb]{0.25,0.63,0.44}{{#1}}}
    \newcommand{\BaseNTok}[1]{\textcolor[rgb]{0.25,0.63,0.44}{{#1}}}
    \newcommand{\FloatTok}[1]{\textcolor[rgb]{0.25,0.63,0.44}{{#1}}}
    \newcommand{\CharTok}[1]{\textcolor[rgb]{0.25,0.44,0.63}{{#1}}}
    \newcommand{\StringTok}[1]{\textcolor[rgb]{0.25,0.44,0.63}{{#1}}}
    \newcommand{\CommentTok}[1]{\textcolor[rgb]{0.38,0.63,0.69}{\textit{{#1}}}}
    \newcommand{\OtherTok}[1]{\textcolor[rgb]{0.00,0.44,0.13}{{#1}}}
    \newcommand{\AlertTok}[1]{\textcolor[rgb]{1.00,0.00,0.00}{\textbf{{#1}}}}
    \newcommand{\FunctionTok}[1]{\textcolor[rgb]{0.02,0.16,0.49}{{#1}}}
    \newcommand{\RegionMarkerTok}[1]{{#1}}
    \newcommand{\ErrorTok}[1]{\textcolor[rgb]{1.00,0.00,0.00}{\textbf{{#1}}}}
    \newcommand{\NormalTok}[1]{{#1}}

    % Additional commands for more recent versions of Pandoc
    \newcommand{\ConstantTok}[1]{\textcolor[rgb]{0.53,0.00,0.00}{{#1}}}
    \newcommand{\SpecialCharTok}[1]{\textcolor[rgb]{0.25,0.44,0.63}{{#1}}}
    \newcommand{\VerbatimStringTok}[1]{\textcolor[rgb]{0.25,0.44,0.63}{{#1}}}
    \newcommand{\SpecialStringTok}[1]{\textcolor[rgb]{0.73,0.40,0.53}{{#1}}}
    \newcommand{\ImportTok}[1]{{#1}}
    \newcommand{\DocumentationTok}[1]{\textcolor[rgb]{0.73,0.13,0.13}{\textit{{#1}}}}
    \newcommand{\AnnotationTok}[1]{\textcolor[rgb]{0.38,0.63,0.69}{\textbf{\textit{{#1}}}}}
    \newcommand{\CommentVarTok}[1]{\textcolor[rgb]{0.38,0.63,0.69}{\textbf{\textit{{#1}}}}}
    \newcommand{\VariableTok}[1]{\textcolor[rgb]{0.10,0.09,0.49}{{#1}}}
    \newcommand{\ControlFlowTok}[1]{\textcolor[rgb]{0.00,0.44,0.13}{\textbf{{#1}}}}
    \newcommand{\OperatorTok}[1]{\textcolor[rgb]{0.40,0.40,0.40}{{#1}}}
    \newcommand{\BuiltInTok}[1]{{#1}}
    \newcommand{\ExtensionTok}[1]{{#1}}
    \newcommand{\PreprocessorTok}[1]{\textcolor[rgb]{0.74,0.48,0.00}{{#1}}}
    \newcommand{\AttributeTok}[1]{\textcolor[rgb]{0.49,0.56,0.16}{{#1}}}
    \newcommand{\InformationTok}[1]{\textcolor[rgb]{0.38,0.63,0.69}{\textbf{\textit{{#1}}}}}
    \newcommand{\WarningTok}[1]{\textcolor[rgb]{0.38,0.63,0.69}{\textbf{\textit{{#1}}}}}
    \makeatletter
    \newsavebox\pandoc@box
    \newcommand*\pandocbounded[1]{%
      \sbox\pandoc@box{#1}%
      % scaling factors for width and height
      \Gscale@div\@tempa\textheight{\dimexpr\ht\pandoc@box+\dp\pandoc@box\relax}%
      \Gscale@div\@tempb\linewidth{\wd\pandoc@box}%
      % select the smaller of both
      \ifdim\@tempb\p@<\@tempa\p@
        \let\@tempa\@tempb
      \fi
      % scaling accordingly (\@tempa < 1)
      \ifdim\@tempa\p@<\p@
        \scalebox{\@tempa}{\usebox\pandoc@box}%
      % scaling not needed, use as it is
      \else
        \usebox{\pandoc@box}%
      \fi
    }
    \makeatother

    % Define a nice break command that doesn't care if a line doesn't already
    % exist.
    \def\br{\hspace*{\fill} \\* }
    % Math Jax compatibility definitions
    \def\gt{>}
    \def\lt{<}
    \let\Oldtex\TeX
    \let\Oldlatex\LaTeX
    \renewcommand{\TeX}{\textrm{\Oldtex}}
    \renewcommand{\LaTeX}{\textrm{\Oldlatex}}
    % Document parameters
    % Document title
    \title{Homework 2\\[0.5ex]\large Computer Simulation 1}
    \author{Andrea Pokraj\v{c}i\'{c}\\
    \normalsize\href{mailto:andrea.pokrajcic@studserv.uni-leipzig.de}{andrea.pokrajcic@studserv.uni-leipzig.de}\\
    \normalsize Matriculation number: 3793956}
    \date{}
    
    
    
    
    
    
    
% Pygments definitions
\makeatletter
\def\PY@reset{\let\PY@it=\relax \let\PY@bf=\relax%
    \let\PY@ul=\relax \let\PY@tc=\relax%
    \let\PY@bc=\relax \let\PY@ff=\relax}
\def\PY@tok#1{\csname PY@tok@#1\endcsname}
\def\PY@toks#1+{\ifx\relax#1\empty\else%
    \PY@tok{#1}\expandafter\PY@toks\fi}
\def\PY@do#1{\PY@bc{\PY@tc{\PY@ul{%
    \PY@it{\PY@bf{\PY@ff{#1}}}}}}}
\def\PY#1#2{\PY@reset\PY@toks#1+\relax+\PY@do{#2}}

\@namedef{PY@tok@w}{\def\PY@tc##1{\textcolor[rgb]{0.73,0.73,0.73}{##1}}}
\@namedef{PY@tok@c}{\let\PY@it=\textit\def\PY@tc##1{\textcolor[rgb]{0.24,0.48,0.48}{##1}}}
\@namedef{PY@tok@cp}{\def\PY@tc##1{\textcolor[rgb]{0.61,0.40,0.00}{##1}}}
\@namedef{PY@tok@k}{\let\PY@bf=\textbf\def\PY@tc##1{\textcolor[rgb]{0.00,0.50,0.00}{##1}}}
\@namedef{PY@tok@kp}{\def\PY@tc##1{\textcolor[rgb]{0.00,0.50,0.00}{##1}}}
\@namedef{PY@tok@kt}{\def\PY@tc##1{\textcolor[rgb]{0.69,0.00,0.25}{##1}}}
\@namedef{PY@tok@o}{\def\PY@tc##1{\textcolor[rgb]{0.40,0.40,0.40}{##1}}}
\@namedef{PY@tok@ow}{\let\PY@bf=\textbf\def\PY@tc##1{\textcolor[rgb]{0.67,0.13,1.00}{##1}}}
\@namedef{PY@tok@nb}{\def\PY@tc##1{\textcolor[rgb]{0.00,0.50,0.00}{##1}}}
\@namedef{PY@tok@nf}{\def\PY@tc##1{\textcolor[rgb]{0.00,0.00,1.00}{##1}}}
\@namedef{PY@tok@nc}{\let\PY@bf=\textbf\def\PY@tc##1{\textcolor[rgb]{0.00,0.00,1.00}{##1}}}
\@namedef{PY@tok@nn}{\let\PY@bf=\textbf\def\PY@tc##1{\textcolor[rgb]{0.00,0.00,1.00}{##1}}}
\@namedef{PY@tok@ne}{\let\PY@bf=\textbf\def\PY@tc##1{\textcolor[rgb]{0.80,0.25,0.22}{##1}}}
\@namedef{PY@tok@nv}{\def\PY@tc##1{\textcolor[rgb]{0.10,0.09,0.49}{##1}}}
\@namedef{PY@tok@no}{\def\PY@tc##1{\textcolor[rgb]{0.53,0.00,0.00}{##1}}}
\@namedef{PY@tok@nl}{\def\PY@tc##1{\textcolor[rgb]{0.46,0.46,0.00}{##1}}}
\@namedef{PY@tok@ni}{\let\PY@bf=\textbf\def\PY@tc##1{\textcolor[rgb]{0.44,0.44,0.44}{##1}}}
\@namedef{PY@tok@na}{\def\PY@tc##1{\textcolor[rgb]{0.41,0.47,0.13}{##1}}}
\@namedef{PY@tok@nt}{\let\PY@bf=\textbf\def\PY@tc##1{\textcolor[rgb]{0.00,0.50,0.00}{##1}}}
\@namedef{PY@tok@nd}{\def\PY@tc##1{\textcolor[rgb]{0.67,0.13,1.00}{##1}}}
\@namedef{PY@tok@s}{\def\PY@tc##1{\textcolor[rgb]{0.73,0.13,0.13}{##1}}}
\@namedef{PY@tok@sd}{\let\PY@it=\textit\def\PY@tc##1{\textcolor[rgb]{0.73,0.13,0.13}{##1}}}
\@namedef{PY@tok@si}{\let\PY@bf=\textbf\def\PY@tc##1{\textcolor[rgb]{0.64,0.35,0.47}{##1}}}
\@namedef{PY@tok@se}{\let\PY@bf=\textbf\def\PY@tc##1{\textcolor[rgb]{0.67,0.36,0.12}{##1}}}
\@namedef{PY@tok@sr}{\def\PY@tc##1{\textcolor[rgb]{0.64,0.35,0.47}{##1}}}
\@namedef{PY@tok@ss}{\def\PY@tc##1{\textcolor[rgb]{0.10,0.09,0.49}{##1}}}
\@namedef{PY@tok@sx}{\def\PY@tc##1{\textcolor[rgb]{0.00,0.50,0.00}{##1}}}
\@namedef{PY@tok@m}{\def\PY@tc##1{\textcolor[rgb]{0.40,0.40,0.40}{##1}}}
\@namedef{PY@tok@gh}{\let\PY@bf=\textbf\def\PY@tc##1{\textcolor[rgb]{0.00,0.00,0.50}{##1}}}
\@namedef{PY@tok@gu}{\let\PY@bf=\textbf\def\PY@tc##1{\textcolor[rgb]{0.50,0.00,0.50}{##1}}}
\@namedef{PY@tok@gd}{\def\PY@tc##1{\textcolor[rgb]{0.63,0.00,0.00}{##1}}}
\@namedef{PY@tok@gi}{\def\PY@tc##1{\textcolor[rgb]{0.00,0.52,0.00}{##1}}}
\@namedef{PY@tok@gr}{\def\PY@tc##1{\textcolor[rgb]{0.89,0.00,0.00}{##1}}}
\@namedef{PY@tok@ge}{\let\PY@it=\textit}
\@namedef{PY@tok@gs}{\let\PY@bf=\textbf}
\@namedef{PY@tok@ges}{\let\PY@bf=\textbf\let\PY@it=\textit}
\@namedef{PY@tok@gp}{\let\PY@bf=\textbf\def\PY@tc##1{\textcolor[rgb]{0.00,0.00,0.50}{##1}}}
\@namedef{PY@tok@go}{\def\PY@tc##1{\textcolor[rgb]{0.44,0.44,0.44}{##1}}}
\@namedef{PY@tok@gt}{\def\PY@tc##1{\textcolor[rgb]{0.00,0.27,0.87}{##1}}}
\@namedef{PY@tok@err}{\def\PY@bc##1{{\setlength{\fboxsep}{\string -\fboxrule}\fcolorbox[rgb]{1.00,0.00,0.00}{1,1,1}{\strut ##1}}}}
\@namedef{PY@tok@kc}{\let\PY@bf=\textbf\def\PY@tc##1{\textcolor[rgb]{0.00,0.50,0.00}{##1}}}
\@namedef{PY@tok@kd}{\let\PY@bf=\textbf\def\PY@tc##1{\textcolor[rgb]{0.00,0.50,0.00}{##1}}}
\@namedef{PY@tok@kn}{\let\PY@bf=\textbf\def\PY@tc##1{\textcolor[rgb]{0.00,0.50,0.00}{##1}}}
\@namedef{PY@tok@kr}{\let\PY@bf=\textbf\def\PY@tc##1{\textcolor[rgb]{0.00,0.50,0.00}{##1}}}
\@namedef{PY@tok@bp}{\def\PY@tc##1{\textcolor[rgb]{0.00,0.50,0.00}{##1}}}
\@namedef{PY@tok@fm}{\def\PY@tc##1{\textcolor[rgb]{0.00,0.00,1.00}{##1}}}
\@namedef{PY@tok@vc}{\def\PY@tc##1{\textcolor[rgb]{0.10,0.09,0.49}{##1}}}
\@namedef{PY@tok@vg}{\def\PY@tc##1{\textcolor[rgb]{0.10,0.09,0.49}{##1}}}
\@namedef{PY@tok@vi}{\def\PY@tc##1{\textcolor[rgb]{0.10,0.09,0.49}{##1}}}
\@namedef{PY@tok@vm}{\def\PY@tc##1{\textcolor[rgb]{0.10,0.09,0.49}{##1}}}
\@namedef{PY@tok@sa}{\def\PY@tc##1{\textcolor[rgb]{0.73,0.13,0.13}{##1}}}
\@namedef{PY@tok@sb}{\def\PY@tc##1{\textcolor[rgb]{0.73,0.13,0.13}{##1}}}
\@namedef{PY@tok@sc}{\def\PY@tc##1{\textcolor[rgb]{0.73,0.13,0.13}{##1}}}
\@namedef{PY@tok@dl}{\def\PY@tc##1{\textcolor[rgb]{0.73,0.13,0.13}{##1}}}
\@namedef{PY@tok@s2}{\def\PY@tc##1{\textcolor[rgb]{0.73,0.13,0.13}{##1}}}
\@namedef{PY@tok@sh}{\def\PY@tc##1{\textcolor[rgb]{0.73,0.13,0.13}{##1}}}
\@namedef{PY@tok@s1}{\def\PY@tc##1{\textcolor[rgb]{0.73,0.13,0.13}{##1}}}
\@namedef{PY@tok@mb}{\def\PY@tc##1{\textcolor[rgb]{0.40,0.40,0.40}{##1}}}
\@namedef{PY@tok@mf}{\def\PY@tc##1{\textcolor[rgb]{0.40,0.40,0.40}{##1}}}
\@namedef{PY@tok@mh}{\def\PY@tc##1{\textcolor[rgb]{0.40,0.40,0.40}{##1}}}
\@namedef{PY@tok@mi}{\def\PY@tc##1{\textcolor[rgb]{0.40,0.40,0.40}{##1}}}
\@namedef{PY@tok@il}{\def\PY@tc##1{\textcolor[rgb]{0.40,0.40,0.40}{##1}}}
\@namedef{PY@tok@mo}{\def\PY@tc##1{\textcolor[rgb]{0.40,0.40,0.40}{##1}}}
\@namedef{PY@tok@ch}{\let\PY@it=\textit\def\PY@tc##1{\textcolor[rgb]{0.24,0.48,0.48}{##1}}}
\@namedef{PY@tok@cm}{\let\PY@it=\textit\def\PY@tc##1{\textcolor[rgb]{0.24,0.48,0.48}{##1}}}
\@namedef{PY@tok@cpf}{\let\PY@it=\textit\def\PY@tc##1{\textcolor[rgb]{0.24,0.48,0.48}{##1}}}
\@namedef{PY@tok@c1}{\let\PY@it=\textit\def\PY@tc##1{\textcolor[rgb]{0.24,0.48,0.48}{##1}}}
\@namedef{PY@tok@cs}{\let\PY@it=\textit\def\PY@tc##1{\textcolor[rgb]{0.24,0.48,0.48}{##1}}}

\def\PYZbs{\char`\\}
\def\PYZus{\char`\_}
\def\PYZob{\char`\{}
\def\PYZcb{\char`\}}
\def\PYZca{\char`\^}
\def\PYZam{\char`\&}
\def\PYZlt{\char`\<}
\def\PYZgt{\char`\>}
\def\PYZsh{\char`\#}
\def\PYZpc{\char`\%}
\def\PYZdl{\char`\$}
\def\PYZhy{\char`\-}
\def\PYZsq{\char`\'}
\def\PYZdq{\char`\"}
\def\PYZti{\char`\~}
% for compatibility with earlier versions
\def\PYZat{@}
\def\PYZlb{[}
\def\PYZrb{]}
\makeatother


    % For linebreaks inside Verbatim environment from package fancyvrb.
    \makeatletter
        \newbox\Wrappedcontinuationbox
        \newbox\Wrappedvisiblespacebox
        \newcommand*\Wrappedvisiblespace {\textcolor{red}{\textvisiblespace}}
        \newcommand*\Wrappedcontinuationsymbol {\textcolor{red}{\llap{\tiny$\m@th\hookrightarrow$}}}
        \newcommand*\Wrappedcontinuationindent {3ex }
        \newcommand*\Wrappedafterbreak {\kern\Wrappedcontinuationindent\copy\Wrappedcontinuationbox}
        % Take advantage of the already applied Pygments mark-up to insert
        % potential linebreaks for TeX processing.
        %        {, <, #, %, $, ' and ": go to next line.
        %        _, }, ^, &, >, - and ~: stay at end of broken line.
        % Use of \textquotesingle for straight quote.
        \newcommand*\Wrappedbreaksatspecials {%
            \def\PYGZus{\discretionary{\char`\_}{\Wrappedafterbreak}{\char`\_}}%
            \def\PYGZob{\discretionary{}{\Wrappedafterbreak\char`\{}{\char`\{}}%
            \def\PYGZcb{\discretionary{\char`\}}{\Wrappedafterbreak}{\char`\}}}%
            \def\PYGZca{\discretionary{\char`\^}{\Wrappedafterbreak}{\char`\^}}%
            \def\PYGZam{\discretionary{\char`\&}{\Wrappedafterbreak}{\char`\&}}%
            \def\PYGZlt{\discretionary{}{\Wrappedafterbreak\char`\<}{\char`\<}}%
            \def\PYGZgt{\discretionary{\char`\>}{\Wrappedafterbreak}{\char`\>}}%
            \def\PYGZsh{\discretionary{}{\Wrappedafterbreak\char`\#}{\char`\#}}%
            \def\PYGZpc{\discretionary{}{\Wrappedafterbreak\char`\%}{\char`\%}}%
            \def\PYGZdl{\discretionary{}{\Wrappedafterbreak\char`\$}{\char`\$}}%
            \def\PYGZhy{\discretionary{\char`\-}{\Wrappedafterbreak}{\char`\-}}%
            \def\PYGZsq{\discretionary{}{\Wrappedafterbreak\textquotesingle}{\textquotesingle}}%
            \def\PYGZdq{\discretionary{}{\Wrappedafterbreak\char`\"}{\char`\"}}%
            \def\PYGZti{\discretionary{\char`\~}{\Wrappedafterbreak}{\char`\~}}%
        }
        % Some characters . , ; ? ! / are not pygmentized.
        % This macro makes them "active" and they will insert potential linebreaks
        \newcommand*\Wrappedbreaksatpunct {%
            \lccode`\~`\.\lowercase{\def~}{\discretionary{\hbox{\char`\.}}{\Wrappedafterbreak}{\hbox{\char`\.}}}%
            \lccode`\~`\,\lowercase{\def~}{\discretionary{\hbox{\char`\,}}{\Wrappedafterbreak}{\hbox{\char`\,}}}%
            \lccode`\~`\;\lowercase{\def~}{\discretionary{\hbox{\char`\;}}{\Wrappedafterbreak}{\hbox{\char`\;}}}%
            \lccode`\~`\:\lowercase{\def~}{\discretionary{\hbox{\char`\:}}{\Wrappedafterbreak}{\hbox{\char`\:}}}%
            \lccode`\~`\?\lowercase{\def~}{\discretionary{\hbox{\char`\?}}{\Wrappedafterbreak}{\hbox{\char`\?}}}%
            \lccode`\~`\!\lowercase{\def~}{\discretionary{\hbox{\char`\!}}{\Wrappedafterbreak}{\hbox{\char`\!}}}%
            \lccode`\~`\/\lowercase{\def~}{\discretionary{\hbox{\char`\/}}{\Wrappedafterbreak}{\hbox{\char`\/}}}%
            \catcode`\.\active
            \catcode`\,\active
            \catcode`\;\active
            \catcode`\:\active
            \catcode`\?\active
            \catcode`\!\active
            \catcode`\/\active
            \lccode`\~`\~
        }
    \makeatother

    \let\OriginalVerbatim=\Verbatim
    \makeatletter
    \renewcommand{\Verbatim}[1][1]{%
        %\parskip\z@skip
        \sbox\Wrappedcontinuationbox {\Wrappedcontinuationsymbol}%
        \sbox\Wrappedvisiblespacebox {\FV@SetupFont\Wrappedvisiblespace}%
        \def\FancyVerbFormatLine ##1{\hsize\linewidth
            \vtop{\raggedright\hyphenpenalty\z@\exhyphenpenalty\z@
                \doublehyphendemerits\z@\finalhyphendemerits\z@
                \strut ##1\strut}%
        }%
        % If the linebreak is at a space, the latter will be displayed as visible
        % space at end of first line, and a continuation symbol starts next line.
        % Stretch/shrink are however usually zero for typewriter font.
        \def\FV@Space {%
            \nobreak\hskip\z@ plus\fontdimen3\font minus\fontdimen4\font
            \discretionary{\copy\Wrappedvisiblespacebox}{\Wrappedafterbreak}
            {\kern\fontdimen2\font}%
        }%

        % Allow breaks at special characters using \PYG... macros.
        \Wrappedbreaksatspecials
        % Breaks at punctuation characters . , ; ? ! and / need catcode=\active
        \OriginalVerbatim[#1,codes*=\Wrappedbreaksatpunct]%
    }
    \makeatother

    % Exact colors from NB
    \definecolor{incolor}{HTML}{303F9F}
    \definecolor{outcolor}{HTML}{D84315}
    \definecolor{cellborder}{HTML}{CFCFCF}
    \definecolor{cellbackground}{HTML}{F7F7F7}

    % prompt
    \makeatletter
    \newcommand{\boxspacing}{\kern\kvtcb@left@rule\kern\kvtcb@boxsep}
    \makeatother
    \newcommand{\prompt}[4]{
        {\ttfamily\llap{{\color{#2}[#3]:\hspace{3pt}#4}}\vspace{-\baselineskip}}
    }
    

    
    % Prevent overflowing lines due to hard-to-break entities
    \sloppy
    % Setup hyperref package
    \hypersetup{
      breaklinks=true,  % so long urls are correctly broken across lines
      colorlinks=true,
      urlcolor=urlcolor,
      linkcolor=linkcolor,
      citecolor=citecolor,
      }
    % Slightly bigger margins than the latex defaults
    
    \geometry{verbose,tmargin=1in,bmargin=1in,lmargin=1in,rmargin=1in}
    
    

\begin{document}
    
    \maketitle
    
    

    
    \section*{Problem 3: Statement}\label{problem-3-statement}

\begin{enumerate}
\def\labelenumi{\alph{enumi})}
\item
  For the example of the one-dimensional Ising model with free boundary
  conditions determine the partition function by exact enumeration of
  all states. To this end write a computer program for \(N = 5, 10\) and
  \(20\) spins (\(2^N = 2^{20} \approx 10^6\) states) at \(\beta=1\)
  (resp. \(T = J/k_B\)). Use the simplest possible implementation
  (\(5-20\) nested do-loops). Check the result for \(N = 5\)
  (\(2^5 = 32\) states) through explicit counting ``by hand''. Finally
  compare with the analytical result of problem 1.
\item
  How can the code be modified with minimal changes such that the exact
  result can be easily obtained for all temperatures?
\item
  With another small modification one can obtain exact results for all
  temperatures and all magnetic fields? How? There will be another
  homework dealing with this problem.
\end{enumerate}

    \section*{Problem 3: Solution}\label{problem-3-solution}

    \subsection*{Part (a)}

From Homework 1, Problem 1 we know that for an 1D Ising model with free
boundaries the partition function can be written as:

\[
Z=\sum_{s_1=\pm1}\cdots\sum_{s_N=\pm1}e^{\frac{J}{k_B T}(s_1s_2+s_2s_3+\cdots+s_{N-2}s_{N-1}+s_{N-1}s_N)} \tag{1}
\]

We can algorithmically determine this partition function by the method
of exact enumaration of all the states.

We do so by using nested \texttt{for} loops, where we loop over the spin
states \([-1,1]\) for each particle in the 1D Ising chain. The number of
nested for loops depends on \(N\), so for \(N=5\), we have 5 loops, for
\(N=10\), we have 10 loops and so on. The energy is defined in the
inner-most nested loop; as given by the exponent of the exponential in
eq. (1). Before performing the calculations, we predefine the variables
and define the exponential function.

    \begin{tcolorbox}[breakable, size=fbox, boxrule=1pt, pad at break*=1mm,colback=cellbackground, colframe=cellborder]
\prompt{In}{incolor}{14}{\boxspacing}
\begin{Verbatim}[commandchars=\\\{\}]
\PY{k+kn}{import}\PY{+w}{ }\PY{n+nn}{math}

\PY{k}{def}\PY{+w}{ }\PY{n+nf}{exp}\PY{p}{(}\PY{n}{x}\PY{p}{)}\PY{p}{:}
    \PY{k}{return} \PY{n}{math}\PY{o}{.}\PY{n}{exp}\PY{p}{(}\PY{n}{x}\PY{p}{)}

\PY{n}{spins}\PY{o}{=}\PY{p}{[}\PY{o}{\PYZhy{}}\PY{l+m+mi}{1}\PY{p}{,}\PY{l+m+mi}{1}\PY{p}{]}

\PY{n}{Z\PYZus{}5}\PY{o}{=}\PY{l+m+mi}{0}
\PY{n}{E\PYZus{}5}\PY{o}{=}\PY{l+m+mi}{0}
\PY{n+nb}{print}\PY{p}{(}\PY{l+s+s2}{\PYZdq{}}\PY{l+s+s2}{EXACT ENUMERATION RESULTS:}\PY{l+s+s2}{\PYZdq{}}\PY{p}{)}


\PY{c+c1}{\PYZsh{}\PYZsh{} N=5}
\PY{k}{for} \PY{n}{s1} \PY{o+ow}{in} \PY{n}{spins}\PY{p}{:}
    \PY{k}{for} \PY{n}{s2} \PY{o+ow}{in} \PY{n}{spins}\PY{p}{:}
  \PY{k}{for} \PY{n}{s3} \PY{o+ow}{in} \PY{n}{spins}\PY{p}{:}
   \PY{k}{for} \PY{n}{s4} \PY{o+ow}{in} \PY{n}{spins}\PY{p}{:}
    \PY{k}{for} \PY{n}{s5} \PY{o+ow}{in} \PY{n}{spins}\PY{p}{:}
     \PY{n}{E\PYZus{}5}\PY{o}{=}\PY{o}{\PYZhy{}}\PY{l+m+mi}{1}\PY{o}{*}\PY{p}{(}\PY{n}{s1}\PY{o}{*}\PY{n}{s2} \PY{o}{+} \PY{n}{s2}\PY{o}{*}\PY{n}{s3}\PY{o}{+}\PY{n}{s3}\PY{o}{*}\PY{n}{s4}\PY{o}{+}\PY{n}{s4}\PY{o}{*}\PY{n}{s5}\PY{p}{)}
     \PY{n}{Z\PYZus{}5}\PY{o}{+}\PY{o}{=}\PY{n}{exp}\PY{p}{(}\PY{o}{\PYZhy{}}\PY{n}{E\PYZus{}5}\PY{p}{)}

\PY{n+nb}{print}\PY{p}{(}\PY{l+s+s2}{\PYZdq{}}\PY{l+s+s2}{For N=5, Z =}\PY{l+s+s2}{\PYZdq{}}\PY{p}{,} \PY{n}{Z\PYZus{}5}\PY{p}{)}

\PY{n}{Z\PYZus{}10}\PY{o}{=}\PY{l+m+mi}{0}
\PY{n}{E\PYZus{}10}\PY{o}{=}\PY{l+m+mi}{0}
\PY{k}{for} \PY{n}{s1} \PY{o+ow}{in} \PY{n}{spins}\PY{p}{:}
    \PY{k}{for} \PY{n}{s2} \PY{o+ow}{in} \PY{n}{spins}\PY{p}{:}
  \PY{k}{for} \PY{n}{s3} \PY{o+ow}{in} \PY{n}{spins}\PY{p}{:}
   \PY{k}{for} \PY{n}{s4} \PY{o+ow}{in} \PY{n}{spins}\PY{p}{:}
    \PY{k}{for} \PY{n}{s5} \PY{o+ow}{in} \PY{n}{spins}\PY{p}{:}
     \PY{k}{for} \PY{n}{s6} \PY{o+ow}{in} \PY{n}{spins}\PY{p}{:}
      \PY{k}{for} \PY{n}{s7} \PY{o+ow}{in} \PY{n}{spins}\PY{p}{:}
       \PY{k}{for} \PY{n}{s8} \PY{o+ow}{in} \PY{n}{spins}\PY{p}{:}
        \PY{k}{for} \PY{n}{s9} \PY{o+ow}{in} \PY{n}{spins}\PY{p}{:}
         \PY{k}{for} \PY{n}{s10} \PY{o+ow}{in} \PY{n}{spins}\PY{p}{:}
          \PY{n}{E\PYZus{}10}\PY{o}{=}\PY{o}{\PYZhy{}}\PY{l+m+mi}{1}\PY{o}{*}\PY{p}{(}\PY{n}{s1}\PY{o}{*}\PY{n}{s2} \PY{o}{+} \PY{n}{s2}\PY{o}{*}\PY{n}{s3}\PY{o}{+}\PY{n}{s3}\PY{o}{*}\PY{n}{s4}\PY{o}{+}\PY{n}{s4}\PY{o}{*}\PY{n}{s5}\PY{o}{+}\PY{n}{s5}\PY{o}{*}\PY{n}{s6}\PY{o}{+}\PY{n}{s6}\PY{o}{*}\PY{n}{s7}\PY{o}{+}\PY{n}{s7}\PY{o}{*}\PY{n}{s8}\PY{o}{+}\PY{n}{s8}\PY{o}{*}\PY{n}{s9}\PY{o}{+}\PY{n}{s9}\PY{o}{*}\PY{n}{s10}\PY{p}{)}
          \PY{n}{Z\PYZus{}10}\PY{o}{+}\PY{o}{=}\PY{n}{exp}\PY{p}{(}\PY{o}{\PYZhy{}}\PY{n}{E\PYZus{}10}\PY{p}{)}
\PY{n+nb}{print}\PY{p}{(}\PY{l+s+s2}{\PYZdq{}}\PY{l+s+s2}{For N=10, Z =}\PY{l+s+s2}{\PYZdq{}}\PY{p}{,} \PY{n}{Z\PYZus{}10}\PY{p}{)}

\PY{n}{Z\PYZus{}20}\PY{o}{=}\PY{l+m+mi}{0}
\PY{n}{E\PYZus{}20}\PY{o}{=}\PY{l+m+mi}{0}
\PY{k}{for} \PY{n}{s1} \PY{o+ow}{in} \PY{n}{spins}\PY{p}{:}
    \PY{k}{for} \PY{n}{s2} \PY{o+ow}{in} \PY{n}{spins}\PY{p}{:}
  \PY{k}{for} \PY{n}{s3} \PY{o+ow}{in} \PY{n}{spins}\PY{p}{:}
   \PY{k}{for} \PY{n}{s4} \PY{o+ow}{in} \PY{n}{spins}\PY{p}{:}
    \PY{k}{for} \PY{n}{s5} \PY{o+ow}{in} \PY{n}{spins}\PY{p}{:}
     \PY{k}{for} \PY{n}{s6} \PY{o+ow}{in} \PY{n}{spins}\PY{p}{:}
      \PY{k}{for} \PY{n}{s7} \PY{o+ow}{in} \PY{n}{spins}\PY{p}{:}
       \PY{k}{for} \PY{n}{s8} \PY{o+ow}{in} \PY{n}{spins}\PY{p}{:}
        \PY{k}{for} \PY{n}{s9} \PY{o+ow}{in} \PY{n}{spins}\PY{p}{:}
         \PY{k}{for} \PY{n}{s10} \PY{o+ow}{in} \PY{n}{spins}\PY{p}{:}
          \PY{k}{for} \PY{n}{s11} \PY{o+ow}{in} \PY{n}{spins}\PY{p}{:}
           \PY{k}{for} \PY{n}{s12} \PY{o+ow}{in} \PY{n}{spins}\PY{p}{:}
            \PY{k}{for} \PY{n}{s13} \PY{o+ow}{in} \PY{n}{spins}\PY{p}{:}
             \PY{k}{for} \PY{n}{s14} \PY{o+ow}{in} \PY{n}{spins}\PY{p}{:}
              \PY{k}{for} \PY{n}{s15} \PY{o+ow}{in} \PY{n}{spins}\PY{p}{:}
               \PY{k}{for} \PY{n}{s16} \PY{o+ow}{in} \PY{n}{spins}\PY{p}{:}
                \PY{k}{for} \PY{n}{s17} \PY{o+ow}{in} \PY{n}{spins}\PY{p}{:}
                 \PY{k}{for} \PY{n}{s18} \PY{o+ow}{in} \PY{n}{spins}\PY{p}{:}
                  \PY{k}{for} \PY{n}{s19} \PY{o+ow}{in} \PY{n}{spins}\PY{p}{:}
                   \PY{k}{for} \PY{n}{s20} \PY{o+ow}{in} \PY{n}{spins}\PY{p}{:}
                    \PY{n}{E\PYZus{}20}\PY{o}{=}\PY{o}{\PYZhy{}}\PY{l+m+mi}{1}\PY{o}{*}\PY{p}{(}\PY{n}{s1}\PY{o}{*}\PY{n}{s2} \PY{o}{+} \PY{n}{s2}\PY{o}{*}\PY{n}{s3} \PY{o}{+} \PY{n}{s3}\PY{o}{*}\PY{n}{s4} \PY{o}{+} \PY{n}{s4}\PY{o}{*}\PY{n}{s5} \PY{o}{+} \PY{n}{s5}\PY{o}{*}\PY{n}{s6} \PY{o}{+} \PY{n}{s6}\PY{o}{*}\PY{n}{s7} \PY{o}{+} \PY{n}{s7}\PY{o}{*}\PY{n}{s8} \PY{o}{+} \PY{n}{s8}\PY{o}{*}\PY{n}{s9} \PY{o}{+} \PY{n}{s9}\PY{o}{*}\PY{n}{s10} \PY{o}{+} \PY{n}{s10}\PY{o}{*}\PY{n}{s11} \PY{o}{+} \PY{n}{s11}\PY{o}{*}\PY{n}{s12} \PY{o}{+} \PY{n}{s12}\PY{o}{*}\PY{n}{s13} \PY{o}{+} \PY{n}{s13}\PY{o}{*}\PY{n}{s14} \PY{o}{+} \PY{n}{s14}\PY{o}{*}\PY{n}{s15} \PY{o}{+} \PY{n}{s15}\PY{o}{*}\PY{n}{s16} \PY{o}{+} \PY{n}{s16}\PY{o}{*}\PY{n}{s17} \PY{o}{+} \PY{n}{s17}\PY{o}{*}\PY{n}{s18} \PY{o}{+} \PY{n}{s18}\PY{o}{*}\PY{n}{s19} \PY{o}{+} \PY{n}{s19}\PY{o}{*}\PY{n}{s20}\PY{p}{)}
                    \PY{n}{Z\PYZus{}20}\PY{o}{+}\PY{o}{=}\PY{n}{exp}\PY{p}{(}\PY{o}{\PYZhy{}}\PY{n}{E\PYZus{}20}\PY{p}{)}

\PY{n+nb}{print}\PY{p}{(}\PY{l+s+s2}{\PYZdq{}}\PY{l+s+s2}{For N=20, Z =}\PY{l+s+s2}{\PYZdq{}}\PY{p}{,} \PY{n}{Z\PYZus{}20}\PY{p}{)}
\end{Verbatim}
\end{tcolorbox}

    \begin{Verbatim}[commandchars=\\\{\}]
EXACT ENUMERATION RESULTS:
For N=5, Z = 181.42806240140408
For N=10, Z = 50792.26102565571
For N=20, Z = 3980922408.678104
    \end{Verbatim}

    Derived Homework 1, the analytical expression for the partition function
is the following

\[
\boxed{Z=2^N\left[\cosh\left(\frac{J}{k_B T}\right)\right]^{N-1}} \tag{2}
\]

Using that we can compute analytically the values of \(Z\) for
\(N=5,10,20\).

    \begin{tcolorbox}[breakable, size=fbox, boxrule=1pt, pad at break*=1mm,colback=cellbackground, colframe=cellborder]
\prompt{In}{incolor}{15}{\boxspacing}
\begin{Verbatim}[commandchars=\\\{\}]
\PY{k}{def}\PY{+w}{ }\PY{n+nf}{Z}\PY{p}{(}\PY{n}{N}\PY{p}{)}\PY{p}{:}
    \PY{k}{return} \PY{l+m+mi}{2}\PY{o}{*}\PY{o}{*}\PY{n}{N}\PY{o}{*}\PY{p}{(}\PY{n}{math}\PY{o}{.}\PY{n}{cosh}\PY{p}{(}\PY{l+m+mi}{1}\PY{p}{)}\PY{p}{)}\PY{o}{*}\PY{o}{*}\PY{p}{(}\PY{n}{N}\PY{o}{\PYZhy{}}\PY{l+m+mi}{1}\PY{p}{)}

\PY{n+nb}{print}\PY{p}{(}\PY{l+s+s2}{\PYZdq{}}\PY{l+s+s2}{ANALYTICAL RESULTS:}\PY{l+s+s2}{\PYZdq{}}\PY{p}{)}
\PY{n+nb}{print}\PY{p}{(}\PY{l+s+s2}{\PYZdq{}}\PY{l+s+s2}{For N=5, Z =}\PY{l+s+s2}{\PYZdq{}}\PY{p}{,} \PY{n}{Z}\PY{p}{(}\PY{l+m+mi}{5}\PY{p}{)}\PY{p}{)}
\PY{n+nb}{print}\PY{p}{(}\PY{l+s+s2}{\PYZdq{}}\PY{l+s+s2}{For N=10, Z =}\PY{l+s+s2}{\PYZdq{}}\PY{p}{,} \PY{n}{Z}\PY{p}{(}\PY{l+m+mi}{10}\PY{p}{)}\PY{p}{)}
\PY{n+nb}{print}\PY{p}{(}\PY{l+s+s2}{\PYZdq{}}\PY{l+s+s2}{For N=20, Z =}\PY{l+s+s2}{\PYZdq{}}\PY{p}{,} \PY{n}{Z}\PY{p}{(}\PY{l+m+mi}{20}\PY{p}{)}\PY{p}{)}
\end{Verbatim}
\end{tcolorbox}

    \begin{Verbatim}[commandchars=\\\{\}]
ANALYTICAL RESULTS:
For N=5, Z = 181.42806240140402
For N=10, Z = 50792.26102565565
For N=20, Z = 3980922408.72465
    \end{Verbatim}

    \begin{tcolorbox}[breakable, size=fbox, boxrule=1pt, pad at break*=1mm,colback=cellbackground, colframe=cellborder]
\prompt{In}{incolor}{16}{\boxspacing}
\begin{Verbatim}[commandchars=\\\{\}]
\PY{n}{d\PYZus{}5}\PY{o}{=}\PY{n}{Z}\PY{p}{(}\PY{l+m+mi}{5}\PY{p}{)}\PY{o}{\PYZhy{}}\PY{n}{Z\PYZus{}5}
\PY{n}{d\PYZus{}10}\PY{o}{=}\PY{n}{Z}\PY{p}{(}\PY{l+m+mi}{10}\PY{p}{)}\PY{o}{\PYZhy{}}\PY{n}{Z\PYZus{}10}
\PY{n}{d\PYZus{}20}\PY{o}{=}\PY{n}{Z}\PY{p}{(}\PY{l+m+mi}{20}\PY{p}{)}\PY{o}{\PYZhy{}}\PY{n}{Z\PYZus{}20}

\PY{n+nb}{print}\PY{p}{(}\PY{l+s+s2}{\PYZdq{}}\PY{l+s+s2}{DIFFERENCES BETWEEN EXACT ENUMERATION AND ANALYTICAL RESULTS:}\PY{l+s+s2}{\PYZdq{}}\PY{p}{)}
\PY{n+nb}{print}\PY{p}{(}\PY{l+s+s2}{\PYZdq{}}\PY{l+s+s2}{For N=5, difference =}\PY{l+s+s2}{\PYZdq{}}\PY{p}{,} \PY{n}{d\PYZus{}5}\PY{p}{)}
\PY{n+nb}{print}\PY{p}{(}\PY{l+s+s2}{\PYZdq{}}\PY{l+s+s2}{For N=10, difference =}\PY{l+s+s2}{\PYZdq{}}\PY{p}{,} \PY{n}{d\PYZus{}10}\PY{p}{)}
\PY{n+nb}{print}\PY{p}{(}\PY{l+s+s2}{\PYZdq{}}\PY{l+s+s2}{For N=20, difference =}\PY{l+s+s2}{\PYZdq{}}\PY{p}{,} \PY{n}{d\PYZus{}20}\PY{p}{)}
\end{Verbatim}
\end{tcolorbox}

    \begin{Verbatim}[commandchars=\\\{\}]
DIFFERENCES BETWEEN EXACT ENUMERATION AND ANALYTICAL RESULTS:
For N=5, difference = -5.684341886080802e-14
For N=10, difference = -5.820766091346741e-11
For N=20, difference = 0.046545982360839844
    \end{Verbatim}

    Comparing the two results we see that they are almost identical. By
subtracting them we find that for \(N=5, 10\) differences are extremely
small, but for \(N=20\) difference is significantly larger in magnitude,
but still small relative to the actual numerical value of \(Z\). I
assume the discrepancy arises due to numerical implementations of
\(\cosh\), because the value is most certainly rounded to a digit that
furthermore introduces a certain error into the analytical result.

    \subsection*{Part (b)}

The code can be modified by adding a simple temperature factor (instead
of 1) in the energy term so that the exact result can be easily obtained
for all temperatures. See below:

    \begin{tcolorbox}[breakable, size=fbox, boxrule=1pt, pad at break*=1mm,colback=cellbackground, colframe=cellborder]
\prompt{In}{incolor}{17}{\boxspacing}
\begin{Verbatim}[commandchars=\\\{\}]
\PY{k+kn}{import}\PY{+w}{ }\PY{n+nn}{math}

\PY{k}{def}\PY{+w}{ }\PY{n+nf}{exp}\PY{p}{(}\PY{n}{x}\PY{p}{)}\PY{p}{:}
    \PY{k}{return} \PY{n}{math}\PY{o}{.}\PY{n}{exp}\PY{p}{(}\PY{n}{x}\PY{p}{)}

\PY{n}{spins}\PY{o}{=}\PY{p}{[}\PY{o}{\PYZhy{}}\PY{l+m+mi}{1}\PY{p}{,}\PY{l+m+mi}{1}\PY{p}{]}

\PY{n}{Z\PYZus{}5}\PY{o}{=}\PY{l+m+mi}{0}
\PY{n}{E\PYZus{}5}\PY{o}{=}\PY{l+m+mi}{0}
\PY{n}{T}\PY{o}{=}\PY{l+m+mf}{1.0}
\PY{n}{J}\PY{o}{=}\PY{l+m+mf}{1.0}
\PY{n}{k\PYZus{}b}\PY{o}{=}\PY{l+m+mf}{1.0}
\PY{n+nb}{print}\PY{p}{(}\PY{l+s+s2}{\PYZdq{}}\PY{l+s+s2}{EXACT ENUMERATION RESULTS:}\PY{l+s+s2}{\PYZdq{}}\PY{p}{)}


\PY{c+c1}{\PYZsh{}\PYZsh{} N=5}
\PY{k}{for} \PY{n}{s1} \PY{o+ow}{in} \PY{n}{spins}\PY{p}{:}
    \PY{k}{for} \PY{n}{s2} \PY{o+ow}{in} \PY{n}{spins}\PY{p}{:}
  \PY{k}{for} \PY{n}{s3} \PY{o+ow}{in} \PY{n}{spins}\PY{p}{:}
   \PY{k}{for} \PY{n}{s4} \PY{o+ow}{in} \PY{n}{spins}\PY{p}{:}
    \PY{k}{for} \PY{n}{s5} \PY{o+ow}{in} \PY{n}{spins}\PY{p}{:}
     \PY{n}{E\PYZus{}5}\PY{o}{=}\PY{o}{\PYZhy{}}\PY{n}{J}\PY{o}{/}\PY{p}{(}\PY{n}{k\PYZus{}b}\PY{o}{*}\PY{n}{T}\PY{p}{)}\PY{o}{*}\PY{p}{(}\PY{n}{s1}\PY{o}{*}\PY{n}{s2} \PY{o}{+} \PY{n}{s2}\PY{o}{*}\PY{n}{s3}\PY{o}{+}\PY{n}{s3}\PY{o}{*}\PY{n}{s4}\PY{o}{+}\PY{n}{s4}\PY{o}{*}\PY{n}{s5}\PY{p}{)}
     \PY{n}{Z\PYZus{}5}\PY{o}{+}\PY{o}{=}\PY{n}{exp}\PY{p}{(}\PY{o}{\PYZhy{}}\PY{n}{E\PYZus{}5}\PY{p}{)}

\PY{n+nb}{print}\PY{p}{(}\PY{l+s+s2}{\PYZdq{}}\PY{l+s+s2}{For N=5, Z =}\PY{l+s+s2}{\PYZdq{}}\PY{p}{,} \PY{n}{Z\PYZus{}5}\PY{p}{)}

\PY{n}{Z\PYZus{}10}\PY{o}{=}\PY{l+m+mi}{0}
\PY{n}{E\PYZus{}10}\PY{o}{=}\PY{l+m+mi}{0}
\PY{k}{for} \PY{n}{s1} \PY{o+ow}{in} \PY{n}{spins}\PY{p}{:}
    \PY{k}{for} \PY{n}{s2} \PY{o+ow}{in} \PY{n}{spins}\PY{p}{:}
  \PY{k}{for} \PY{n}{s3} \PY{o+ow}{in} \PY{n}{spins}\PY{p}{:}
   \PY{k}{for} \PY{n}{s4} \PY{o+ow}{in} \PY{n}{spins}\PY{p}{:}
    \PY{k}{for} \PY{n}{s5} \PY{o+ow}{in} \PY{n}{spins}\PY{p}{:}
     \PY{k}{for} \PY{n}{s6} \PY{o+ow}{in} \PY{n}{spins}\PY{p}{:}
      \PY{k}{for} \PY{n}{s7} \PY{o+ow}{in} \PY{n}{spins}\PY{p}{:}
       \PY{k}{for} \PY{n}{s8} \PY{o+ow}{in} \PY{n}{spins}\PY{p}{:}
        \PY{k}{for} \PY{n}{s9} \PY{o+ow}{in} \PY{n}{spins}\PY{p}{:}
         \PY{k}{for} \PY{n}{s10} \PY{o+ow}{in} \PY{n}{spins}\PY{p}{:}
          \PY{n}{E\PYZus{}10}\PY{o}{=}\PY{o}{\PYZhy{}}\PY{n}{J}\PY{o}{/}\PY{p}{(}\PY{n}{k\PYZus{}b}\PY{o}{*}\PY{n}{T}\PY{p}{)}\PY{o}{*}\PY{p}{(}\PY{n}{s1}\PY{o}{*}\PY{n}{s2} \PY{o}{+} \PY{n}{s2}\PY{o}{*}\PY{n}{s3}\PY{o}{+}\PY{n}{s3}\PY{o}{*}\PY{n}{s4}\PY{o}{+}\PY{n}{s4}\PY{o}{*}\PY{n}{s5}\PY{o}{+}\PY{n}{s5}\PY{o}{*}\PY{n}{s6}\PY{o}{+}\PY{n}{s6}\PY{o}{*}\PY{n}{s7}\PY{o}{+}\PY{n}{s7}\PY{o}{*}\PY{n}{s8}\PY{o}{+}\PY{n}{s8}\PY{o}{*}\PY{n}{s9}\PY{o}{+}\PY{n}{s9}\PY{o}{*}\PY{n}{s10}\PY{p}{)}
          \PY{n}{Z\PYZus{}10}\PY{o}{+}\PY{o}{=}\PY{n}{exp}\PY{p}{(}\PY{o}{\PYZhy{}}\PY{n}{E\PYZus{}10}\PY{p}{)}
\PY{n+nb}{print}\PY{p}{(}\PY{l+s+s2}{\PYZdq{}}\PY{l+s+s2}{For N=10, Z =}\PY{l+s+s2}{\PYZdq{}}\PY{p}{,} \PY{n}{Z\PYZus{}10}\PY{p}{)}

\PY{n}{Z\PYZus{}20}\PY{o}{=}\PY{l+m+mi}{0}
\PY{n}{E\PYZus{}20}\PY{o}{=}\PY{l+m+mi}{0}
\PY{k}{for} \PY{n}{s1} \PY{o+ow}{in} \PY{n}{spins}\PY{p}{:}
    \PY{k}{for} \PY{n}{s2} \PY{o+ow}{in} \PY{n}{spins}\PY{p}{:}
  \PY{k}{for} \PY{n}{s3} \PY{o+ow}{in} \PY{n}{spins}\PY{p}{:}
   \PY{k}{for} \PY{n}{s4} \PY{o+ow}{in} \PY{n}{spins}\PY{p}{:}
    \PY{k}{for} \PY{n}{s5} \PY{o+ow}{in} \PY{n}{spins}\PY{p}{:}
     \PY{k}{for} \PY{n}{s6} \PY{o+ow}{in} \PY{n}{spins}\PY{p}{:}
      \PY{k}{for} \PY{n}{s7} \PY{o+ow}{in} \PY{n}{spins}\PY{p}{:}
       \PY{k}{for} \PY{n}{s8} \PY{o+ow}{in} \PY{n}{spins}\PY{p}{:}
        \PY{k}{for} \PY{n}{s9} \PY{o+ow}{in} \PY{n}{spins}\PY{p}{:}
         \PY{k}{for} \PY{n}{s10} \PY{o+ow}{in} \PY{n}{spins}\PY{p}{:}
          \PY{k}{for} \PY{n}{s11} \PY{o+ow}{in} \PY{n}{spins}\PY{p}{:}
           \PY{k}{for} \PY{n}{s12} \PY{o+ow}{in} \PY{n}{spins}\PY{p}{:}
            \PY{k}{for} \PY{n}{s13} \PY{o+ow}{in} \PY{n}{spins}\PY{p}{:}
             \PY{k}{for} \PY{n}{s14} \PY{o+ow}{in} \PY{n}{spins}\PY{p}{:}
              \PY{k}{for} \PY{n}{s15} \PY{o+ow}{in} \PY{n}{spins}\PY{p}{:}
               \PY{k}{for} \PY{n}{s16} \PY{o+ow}{in} \PY{n}{spins}\PY{p}{:}
                \PY{k}{for} \PY{n}{s17} \PY{o+ow}{in} \PY{n}{spins}\PY{p}{:}
                 \PY{k}{for} \PY{n}{s18} \PY{o+ow}{in} \PY{n}{spins}\PY{p}{:}
                  \PY{k}{for} \PY{n}{s19} \PY{o+ow}{in} \PY{n}{spins}\PY{p}{:}
                   \PY{k}{for} \PY{n}{s20} \PY{o+ow}{in} \PY{n}{spins}\PY{p}{:}
                    \PY{n}{E\PYZus{}20}\PY{o}{=}\PY{o}{\PYZhy{}}\PY{n}{J}\PY{o}{/}\PY{p}{(}\PY{n}{k\PYZus{}b}\PY{o}{*}\PY{n}{T}\PY{p}{)}\PY{o}{*}\PY{p}{(}\PY{n}{s1}\PY{o}{*}\PY{n}{s2} \PY{o}{+} \PY{n}{s2}\PY{o}{*}\PY{n}{s3} \PY{o}{+} \PY{n}{s3}\PY{o}{*}\PY{n}{s4} \PY{o}{+} \PY{n}{s4}\PY{o}{*}\PY{n}{s5} \PY{o}{+} \PY{n}{s5}\PY{o}{*}\PY{n}{s6} \PY{o}{+} \PY{n}{s6}\PY{o}{*}\PY{n}{s7} \PY{o}{+} \PY{n}{s7}\PY{o}{*}\PY{n}{s8} \PY{o}{+} \PY{n}{s8}\PY{o}{*}\PY{n}{s9} \PY{o}{+} \PY{n}{s9}\PY{o}{*}\PY{n}{s10} \PY{o}{+} \PY{n}{s10}\PY{o}{*}\PY{n}{s11} \PY{o}{+} \PY{n}{s11}\PY{o}{*}\PY{n}{s12} \PY{o}{+} \PY{n}{s12}\PY{o}{*}\PY{n}{s13} \PY{o}{+} \PY{n}{s13}\PY{o}{*}\PY{n}{s14} \PY{o}{+} \PY{n}{s14}\PY{o}{*}\PY{n}{s15} \PY{o}{+} \PY{n}{s15}\PY{o}{*}\PY{n}{s16} \PY{o}{+} \PY{n}{s16}\PY{o}{*}\PY{n}{s17} \PY{o}{+} \PY{n}{s17}\PY{o}{*}\PY{n}{s18} \PY{o}{+} \PY{n}{s18}\PY{o}{*}\PY{n}{s19} \PY{o}{+} \PY{n}{s19}\PY{o}{*}\PY{n}{s20}\PY{p}{)}
                    \PY{n}{Z\PYZus{}20}\PY{o}{+}\PY{o}{=}\PY{n}{exp}\PY{p}{(}\PY{o}{\PYZhy{}}\PY{n}{E\PYZus{}20}\PY{p}{)}

\PY{n+nb}{print}\PY{p}{(}\PY{l+s+s2}{\PYZdq{}}\PY{l+s+s2}{For N=20, Z =}\PY{l+s+s2}{\PYZdq{}}\PY{p}{,} \PY{n}{Z\PYZus{}20}\PY{p}{)}
\end{Verbatim}
\end{tcolorbox}

    \begin{Verbatim}[commandchars=\\\{\}]
EXACT ENUMERATION RESULTS:
For N=5, Z = 181.42806240140408
For N=10, Z = 50792.26102565571
For N=20, Z = 3980922408.678104
    \end{Verbatim}

    \subsection*{Part (c)}

By adding another term into the energy function, we can include a
magnetic field term. The final energy expression is therefore then given
by eq. (3)

\[
E=-J \sum_{i=1}^{N-1}s_i s_{i+1} + B\sum_{i=1}^{N}s_i \tag{3}
\]

This together with the previous temperature accounting change allows us
to calculate the partition function for any temperature and magnetic
field.

    \section*{Problem 4: Statement}\label{problem-4-statement}

The free energy per spin, \(\beta f_N \equiv -\frac{1}{N}\ln(Z_N)\), of
the one-dimensional Ising model of a chain with \(N\) spins reads in the
thermodynamic limit \(N \to \infty\) independent of the boundary
conditions:

\[
\beta f_{\infty} \equiv -\lim_{N \to \infty} \frac{1}{N}\ln(Z_N)=-\ln(2)-\ln(\cosh(\beta))
\]

\textbf{Note: This equation in the problem statement assumes \(J=1\),
hence in my solutions I will too assume \(J=1\)}

\begin{enumerate}
\def\labelenumi{\alph{enumi})}
\item
  Calculate the finite-size correction \(\beta f_N - \beta f_{\infty}\)
  for i) free and ii) periodic boundary conditions.
\item
  Plot both corrections as a function of \(N\) for \(\beta = 1, 2, 10\)
  in the range \([1, 10]\). To be able to disentangle the different
  curves, you may set the y-range to, e.g., \([0, 2]\). For a comparison
  of the two boundary conditions it is useful to plot
  \(|\beta f_N - \beta f_{\infty}|\) (why?).
\item
  Repeat b) with a logarithmic y-scale in the \(N\)-range \([1, 100]\).
  Why does the curve for \(\beta = 1\) and periodic boundary conditions
  appear as a straight line in the limit of large \(N\)? What does that
  mean mathematically? How can one understand the analogous curves for
  \(\beta = 2, 10\)?
\end{enumerate}

    \section*{Problem 4: Solution}\label{problem-4-solution}

    \subsection*{Part (a)}

In this task we wish to calculate the finite-size correction
\(\beta f_N - \beta f_{\infty}\) for both free and periodic boundary
conditions. Let us start with the free boundary condition.

(i) Free boundary conditions

Recalling from Homework 1, the partition function of a 1D Ising model
with free boundary conditions is

\[
\boxed{Z_N=2^N\left[\cosh(\beta)\right]^{N-1}} \tag{4}
\]

From the problem statement, the free energy per spin is

\[
\beta f_N \equiv -\frac{1}{N}\ln(Z_N) \tag{5}
\]

Plugging eq. (4) into eq. (5)

\[
\begin{aligned}
\beta f_N &= -\frac{1}{N}\ln\left(2^N\left[\cosh(\beta)\right]^{N-1}\right)\\
&= -\frac{1}{N}\Big(N\ln(2)+(N-1)\ln(\cosh(\beta))\Big)\\
&= -\ln(2)-\frac{N-1}{N}\ln(\cosh(\beta))
\end{aligned} \tag{6}
\]

Again, from the problem statement we know

\[
\beta f_{\infty} \equiv -\lim_{N \to \infty} \frac{1}{N}\ln(Z_N)=-\ln(2)-\ln(\cosh(\beta)) \tag{7}
\]

Therefore, we can calculate \(\beta f_N - \beta f_{\infty}\) by
subtracting eq. (7) from eq. (6):

\[
\begin{aligned}
\beta f_N - \beta f_{\infty} &= -\ln(2)-\frac{N-1}{N}\ln(\cosh(\beta))+\ln(2)+\ln(\cosh(\beta))\\
&= \left(1-\frac{N-1}{N}\right)\ln(\cosh(\beta))\\
&= \boxed{\frac{1}{N}\ln(\cosh(\beta))}
\end{aligned} \tag{8}
\]

(ii) Periodic boundary conditions

During the lecture we derived the partition function of the 1D Ising
model with periodic boundary conditions using the transfer matrix
method. The derived expression, with \(J=1\), is

\[
\begin{aligned}
Z_N &= 2^N\cosh^N(\beta)\left[1+\tanh^N(\beta)\right]\\
&= 2^N\cosh^N(\beta)+2^N\sinh^N(\beta)
\end{aligned} \tag{9}
\]

Using eq. (5) as a blueprint, the free energy per spin is

\[
\begin{aligned}
\beta f_N &= -\frac{1}{N}\ln\left(2^N\cosh^N(\beta)\left[1+\tanh^N(\beta)\right]\right)\\
&= -\frac{1}{N}\Big(\ln\left(2^N\cosh^N(\beta)\right)+\ln\left(1+\tanh^N(\beta)\right)\Big)\\
&= -\frac{1}{N}\Big(N\ln(2)+N\ln(\cosh(\beta))+\ln\left(1+\tanh^N(\beta)\right)\Big)\\
&= -\ln(2)-\ln(\cosh(\beta))-\frac{1}{N}\ln\left(1+\tanh^N(\beta)\right)
\end{aligned} \tag{10}
\]

Then we calculate \(\beta f_N - \beta f_{\infty}\) by subtracting eq.
(7) from eq. (10):

\[
\begin{aligned}
\beta f_N - \beta f_{\infty} &= -\ln(2)-\ln(\cosh(\beta))-\frac{1}{N}\ln\left(1+\tanh^N(\beta)\right)+\ln(2)+\ln(\cosh(\beta))\\
&= \boxed{-\frac{1}{N}\ln\left(1+\tanh^N(\beta)\right)}
\end{aligned} \tag{11}
\]

    \subsection*{Part (b)}

In this part of the problem we plot both the absolute values of both
corrections a function of \(N\) for \(\beta = 1, 2, 10\) in the range
\([1, 10]\). It is useful to plot the absolute value
\(|\beta f_N - \beta f_{\infty}|\) because the periodic boundary
correction function is always general negative. In order to disentangle
the different curves, we set the y-range to \([0, 2]\).

    \begin{tcolorbox}[breakable, size=fbox, boxrule=1pt, pad at break*=1mm,colback=cellbackground, colframe=cellborder]
\prompt{In}{incolor}{5}{\boxspacing}
\begin{Verbatim}[commandchars=\\\{\}]
\PY{k+kn}{import}\PY{+w}{ }\PY{n+nn}{numpy}\PY{+w}{ }\PY{k}{as}\PY{+w}{ }\PY{n+nn}{np}
\PY{k+kn}{from}\PY{+w}{ }\PY{n+nn}{matplotlib}\PY{+w}{ }\PY{k+kn}{import} \PY{n}{pyplot} \PY{k}{as} \PY{n}{plt}

\PY{n}{J} \PY{o}{=} \PY{l+m+mf}{1.0}
\PY{n}{k\PYZus{}b} \PY{o}{=} \PY{l+m+mf}{1.0}

\PY{n}{beta\PYZus{}values} \PY{o}{=} \PY{p}{[}\PY{l+m+mi}{1}\PY{p}{,}\PY{l+m+mi}{2}\PY{p}{,}\PY{l+m+mi}{10}\PY{p}{]}

\PY{n}{N} \PY{o}{=} \PY{n}{np}\PY{o}{.}\PY{n}{linspace}\PY{p}{(}\PY{l+m+mi}{1}\PY{p}{,} \PY{l+m+mi}{10}\PY{p}{,} \PY{l+m+mi}{500}\PY{p}{)}

\PY{k}{def}\PY{+w}{ }\PY{n+nf}{corr\PYZus{}free}\PY{p}{(}\PY{n}{N}\PY{p}{,}\PY{n}{beta}\PY{p}{)}\PY{p}{:}
   \PY{k}{return} \PY{n+nb}{abs}\PY{p}{(}\PY{l+m+mi}{1}\PY{o}{/}\PY{n}{N}\PY{o}{*}\PY{n}{np}\PY{o}{.}\PY{n}{log}\PY{p}{(}\PY{n}{np}\PY{o}{.}\PY{n}{cosh}\PY{p}{(}\PY{n}{beta}\PY{p}{)}\PY{p}{)}\PY{p}{)}

\PY{k}{def}\PY{+w}{ }\PY{n+nf}{corr\PYZus{}periodic}\PY{p}{(}\PY{n}{N}\PY{p}{,}\PY{n}{beta}\PY{p}{)}\PY{p}{:}
   \PY{k}{return} \PY{n+nb}{abs}\PY{p}{(}\PY{o}{\PYZhy{}}\PY{l+m+mi}{1}\PY{o}{/}\PY{n}{N}\PY{o}{*}\PY{n}{np}\PY{o}{.}\PY{n}{log}\PY{p}{(}\PY{l+m+mi}{1}\PY{o}{+}\PY{n}{np}\PY{o}{.}\PY{n}{tanh}\PY{p}{(}\PY{n}{beta}\PY{p}{)}\PY{o}{*}\PY{o}{*}\PY{n}{N}\PY{p}{)}\PY{p}{)}


\PY{k}{for} \PY{n}{beta} \PY{o+ow}{in} \PY{n}{beta\PYZus{}values}\PY{p}{:}
    \PY{n}{plt}\PY{o}{.}\PY{n}{plot}\PY{p}{(}\PY{n}{N}\PY{p}{,} \PY{n}{corr\PYZus{}free}\PY{p}{(}\PY{n}{N}\PY{p}{,} \PY{n}{beta}\PY{p}{)}\PY{p}{,} \PY{n}{label}\PY{o}{=}\PY{l+s+sa}{f}\PY{l+s+s2}{\PYZdq{}}\PY{l+s+s2}{Free, β=}\PY{l+s+si}{\PYZob{}}\PY{n}{beta}\PY{l+s+si}{\PYZcb{}}\PY{l+s+s2}{\PYZdq{}}\PY{p}{)}
    \PY{n}{plt}\PY{o}{.}\PY{n}{plot}\PY{p}{(}\PY{n}{N}\PY{p}{,} \PY{n}{corr\PYZus{}periodic}\PY{p}{(}\PY{n}{N}\PY{p}{,} \PY{n}{beta}\PY{p}{)}\PY{p}{,} \PY{n}{label}\PY{o}{=}\PY{l+s+sa}{f}\PY{l+s+s2}{\PYZdq{}}\PY{l+s+s2}{Periodic, β=}\PY{l+s+si}{\PYZob{}}\PY{n}{beta}\PY{l+s+si}{\PYZcb{}}\PY{l+s+s2}{\PYZdq{}}\PY{p}{)}

\PY{n}{plt}\PY{o}{.}\PY{n}{ylim}\PY{p}{(}\PY{l+m+mi}{0}\PY{p}{,}\PY{l+m+mi}{2}\PY{p}{)}
\PY{n}{plt}\PY{o}{.}\PY{n}{xlabel}\PY{p}{(}\PY{l+s+s2}{\PYZdq{}}\PY{l+s+s2}{N}\PY{l+s+s2}{\PYZdq{}}\PY{p}{)}
\PY{n}{plt}\PY{o}{.}\PY{n}{ylabel}\PY{p}{(}\PY{l+s+sa}{r}\PY{l+s+s2}{\PYZdq{}}\PY{l+s+s2}{\PYZdl{}}\PY{l+s+s2}{\PYZbs{}}\PY{l+s+s2}{beta f\PYZus{}N \PYZhy{} }\PY{l+s+s2}{\PYZbs{}}\PY{l+s+s2}{beta f\PYZus{}}\PY{l+s+s2}{\PYZob{}}\PY{l+s+s2}{\PYZbs{}}\PY{l+s+s2}{infty\PYZcb{}\PYZdl{}}\PY{l+s+s2}{\PYZdq{}}\PY{p}{)}
\PY{n}{plt}\PY{o}{.}\PY{n}{legend}\PY{p}{(}\PY{p}{)}
\PY{n}{plt}\PY{o}{.}\PY{n}{show}\PY{p}{(}\PY{p}{)}
\end{Verbatim}
\end{tcolorbox}

    \begin{center}
    \adjustimage{max size={0.9\linewidth}{0.9\paperheight}}{CompSim_HW2_files/CompSim_HW2_15_0.png}
    \end{center}
    { \hspace*{\fill} \\}
    
    Figure 1. Dependence of \(|\beta f_N - \beta f_{\infty}|\) on \(N\), for
free and periodic boundary conditions and \(\beta=1,2,10\)

    \subsection*{Part (c)}

In this part of the problem we repeat (b) with a logarithmic y-scale in
the range \(N \in [1, 100]\).

    \begin{tcolorbox}[breakable, size=fbox, boxrule=1pt, pad at break*=1mm,colback=cellbackground, colframe=cellborder]
\prompt{In}{incolor}{7}{\boxspacing}
\begin{Verbatim}[commandchars=\\\{\}]
\PY{k+kn}{import}\PY{+w}{ }\PY{n+nn}{numpy}\PY{+w}{ }\PY{k}{as}\PY{+w}{ }\PY{n+nn}{np}
\PY{k+kn}{from}\PY{+w}{ }\PY{n+nn}{matplotlib}\PY{+w}{ }\PY{k+kn}{import} \PY{n}{pyplot} \PY{k}{as} \PY{n}{plt}

\PY{n}{J} \PY{o}{=} \PY{l+m+mf}{1.0}
\PY{n}{k\PYZus{}b} \PY{o}{=} \PY{l+m+mf}{1.0}

\PY{n}{beta\PYZus{}values} \PY{o}{=} \PY{p}{[}\PY{l+m+mi}{1}\PY{p}{,}\PY{l+m+mi}{2}\PY{p}{,}\PY{l+m+mi}{10}\PY{p}{]}

\PY{n}{N} \PY{o}{=} \PY{n}{np}\PY{o}{.}\PY{n}{linspace}\PY{p}{(}\PY{l+m+mi}{1}\PY{p}{,} \PY{l+m+mi}{100}\PY{p}{,} \PY{l+m+mi}{500}\PY{p}{)}

\PY{k}{def}\PY{+w}{ }\PY{n+nf}{corr\PYZus{}free}\PY{p}{(}\PY{n}{N}\PY{p}{,}\PY{n}{beta}\PY{p}{)}\PY{p}{:}
   \PY{k}{return} \PY{n+nb}{abs}\PY{p}{(}\PY{l+m+mi}{1}\PY{o}{/}\PY{n}{N}\PY{o}{*}\PY{n}{np}\PY{o}{.}\PY{n}{log}\PY{p}{(}\PY{n}{np}\PY{o}{.}\PY{n}{cosh}\PY{p}{(}\PY{n}{beta}\PY{p}{)}\PY{p}{)}\PY{p}{)}

\PY{k}{def}\PY{+w}{ }\PY{n+nf}{corr\PYZus{}periodic}\PY{p}{(}\PY{n}{N}\PY{p}{,}\PY{n}{beta}\PY{p}{)}\PY{p}{:}
   \PY{k}{return} \PY{n+nb}{abs}\PY{p}{(}\PY{o}{\PYZhy{}}\PY{l+m+mi}{1}\PY{o}{/}\PY{n}{N}\PY{o}{*}\PY{n}{np}\PY{o}{.}\PY{n}{log}\PY{p}{(}\PY{l+m+mi}{1}\PY{o}{+}\PY{n}{np}\PY{o}{.}\PY{n}{tanh}\PY{p}{(}\PY{n}{beta}\PY{p}{)}\PY{o}{*}\PY{o}{*}\PY{n}{N}\PY{p}{)}\PY{p}{)}


\PY{k}{for} \PY{n}{beta} \PY{o+ow}{in} \PY{n}{beta\PYZus{}values}\PY{p}{:}
    \PY{n}{plt}\PY{o}{.}\PY{n}{plot}\PY{p}{(}\PY{n}{N}\PY{p}{,} \PY{n}{corr\PYZus{}free}\PY{p}{(}\PY{n}{N}\PY{p}{,} \PY{n}{beta}\PY{p}{)}\PY{p}{,} \PY{n}{label}\PY{o}{=}\PY{l+s+sa}{f}\PY{l+s+s2}{\PYZdq{}}\PY{l+s+s2}{Free, β=}\PY{l+s+si}{\PYZob{}}\PY{n}{beta}\PY{l+s+si}{\PYZcb{}}\PY{l+s+s2}{\PYZdq{}}\PY{p}{)}
    \PY{n}{plt}\PY{o}{.}\PY{n}{plot}\PY{p}{(}\PY{n}{N}\PY{p}{,} \PY{n}{corr\PYZus{}periodic}\PY{p}{(}\PY{n}{N}\PY{p}{,} \PY{n}{beta}\PY{p}{)}\PY{p}{,} \PY{n}{label}\PY{o}{=}\PY{l+s+sa}{f}\PY{l+s+s2}{\PYZdq{}}\PY{l+s+s2}{Periodic, β=}\PY{l+s+si}{\PYZob{}}\PY{n}{beta}\PY{l+s+si}{\PYZcb{}}\PY{l+s+s2}{\PYZdq{}}\PY{p}{)}

\PY{n}{plt}\PY{o}{.}\PY{n}{yscale}\PY{p}{(}\PY{l+s+s2}{\PYZdq{}}\PY{l+s+s2}{log}\PY{l+s+s2}{\PYZdq{}}\PY{p}{)}
\PY{n}{plt}\PY{o}{.}\PY{n}{xlabel}\PY{p}{(}\PY{l+s+s2}{\PYZdq{}}\PY{l+s+s2}{N}\PY{l+s+s2}{\PYZdq{}}\PY{p}{)}
\PY{n}{plt}\PY{o}{.}\PY{n}{ylabel}\PY{p}{(}\PY{l+s+sa}{r}\PY{l+s+s2}{\PYZdq{}}\PY{l+s+s2}{\PYZdl{}}\PY{l+s+s2}{\PYZbs{}}\PY{l+s+s2}{beta f\PYZus{}N \PYZhy{} }\PY{l+s+s2}{\PYZbs{}}\PY{l+s+s2}{beta f\PYZus{}}\PY{l+s+s2}{\PYZob{}}\PY{l+s+s2}{\PYZbs{}}\PY{l+s+s2}{infty\PYZcb{}\PYZdl{}}\PY{l+s+s2}{\PYZdq{}}\PY{p}{)}
\PY{n}{plt}\PY{o}{.}\PY{n}{legend}\PY{p}{(}\PY{p}{)}
\PY{n}{plt}\PY{o}{.}\PY{n}{show}\PY{p}{(}\PY{p}{)}
\end{Verbatim}
\end{tcolorbox}

    \begin{center}
    \adjustimage{max size={0.9\linewidth}{0.9\paperheight}}{CompSim_HW2_files/CompSim_HW2_18_0.png}
    \end{center}
    { \hspace*{\fill} \\}
    
    Figure 2. Dependence of \(|\beta f_N - \beta f_{\infty}|\) on \(N\) on a
logarithmic scale, for free and periodic boundary conditions and
\(\beta=1,2,10\)

    As we can see on Fig. 2, for periodic boundary conditions with
\(\beta = 1\) the curve decreases linearly for large \(N\). Since the
y-axis is logarithmic, this means that on a normal scale the correction
decays exponentially with \(N\).

This can also be seen analytically. Recalling eq. (11), the finite-size
correction for periodic boundary conditions is

\[
\beta f_N - \beta f_{\infty} = -\frac{1}{N}\ln\left(1+\tanh^N(\beta)\right) \tag{12}
\]

Since \(\tanh(1) < 1\), for large \(N\) we have \(\tanh^N(1) \to 0\).
For small \(\tanh^N(1)\) we can therefore use the Taylor expansion
\(\ln(1+y) \approx y\) for \(|y| \ll 1\), which gives

\[
\ln\left(1+\tanh^N(1)\right) \approx \tanh^N(1) \tag{13}
\]

Plugging eq. (13) into eq. (12), for large \(N\) we obtain

\[
\boxed{|\beta f_N - \beta f_{\infty}| \approx \frac{1}{N}\tanh^N(1)} \tag{14}
\]

which is exactly the exponential behaviour observed on Fig. 2. To see
why it appears as a straight line, we take the logarithm of eq. (14):

\[
\begin{aligned}
\ln|\beta f_N - \beta f_{\infty}| &\approx \ln\left(\frac{1}{N}\tanh^N(1)\right)\\
&= N\ln(\tanh(1)) - \ln(N)
\end{aligned} \tag{15}
\]

where \(\ln(\tanh(1)) \approx -0.2723\). For large \(N\) the term
\(N\ln(\tanh(1))\) dominates, since it grows linearly in \(N\), while
\(\ln(N)\) grows only logarithmically. This allows us to neglect the
\(\ln(N)\) term, so the curve on the logarithmic scale is
(approximately) a straight line with slope

\[
\boxed{\ln(\tanh(1)) \approx -0.2723} \tag{16}
\]

Mathematically this means that the finite-size correction decays
exponentially with \(N\), since \(\tanh^N(1) = e^{N\ln(\tanh(1))}\) with
\(\ln(\tanh(1)) < 0\).


    % Add a bibliography block to the postdoc


    % Add a bibliography block to the postdoc
    
    
    
\end{document}
